\begin{description}
\item[Q1] \textbf{Why does the empirical arg-min $r$ match $\lceil N^{2/3}\rceil$ \emph{exactly} at 4/4 tested $N$, when the paper's proof only guarantees the constant to be tight up to logarithmic factors?}\\
\emph{Basis:} at $N\in\{6,8,10,12\}$ the empirical arg-min is $4,4,5,6$, all identically equal to $\lceil N^{2/3}\rceil$; no $\log$-factor drift is visible.\\
\emph{Next steps:} extend the sweep to $N\in\{14,18,24\}$ using sparse-matrix Krylov iteration (avoiding the dense $\binom{2N}{r}$ blowup) to test whether the exact identity persists or drifts by $\pm 1$.

\item[Q2] \textbf{Do the additive $r$ setup cost and the Grover iteration cost cross over at the predicted $r=N^{2/3}$, or does the crossover shift for the ``asymmetric'' regime $M\gg N$?}\\
\emph{Basis:} our simulation used $N=M$; the paper's Theorem~8 explicitly changes shape at $M=N^2$.\\
\emph{Next steps:} port the Johnson-graph enumeration to the two-graph categorical product $J(N,\ell)\times J(M,m)$, sweep $(\ell,m)$ for $M\in\{N,\,N^{3/2},\,N^2\}$, and check whether the empirical arg-min matches the paper's $\ell=\Theta((N^2/M)^{1/3})$ prediction.

\item[Q3] \textbf{Is the coarsened Szegedy walk $U=R_S R_M$ we simulated quantitatively equivalent to the full Szegedy walk on the bipartite doubling, at these small $N$, or does it under-count by an $O(1)$ factor?}\\
\emph{Basis:} the two-dimensional (marked-vs-unmarked) reduction is only \emph{asymptotically} tight; at $N=6$ the empirical $k^\ast=2$ vs $\pi/(4\sqrt{\varepsilon})=2.6$ suggests a small integer-round bias.\\
\emph{Next steps:} explicitly construct the $2\binom{2N}{r}$-dim Szegedy walk unitary and compare $k^\ast$ head-to-head with the coarsened version at $N=6,8,10$; report the discrepancy as a fraction.

\item[Q4] \textbf{How does the Grover-style success amplitude degrade if the planted claw is not unique --- i.e.\ how many extra iterations does the walk need for $|M|>1$, given that the marked-fraction $\varepsilon$ increases linearly in the claw count?}\\
\emph{Basis:} we planted \emph{exactly} one claw; real cryptanalytic instances often have $\Theta(1)$ or $\omega(1)$ claws (e.g.\ birthday collisions in a $2$-to-$1$ function).\\
\emph{Next steps:} rerun the sim with $c\in\{1,2,4,8\}$ planted claws at fixed $N=10$; verify $k^\ast$ scales as $1/\sqrt{c}$ (theory) and that the peak amplitude is unaffected.

\item[Q5] \textbf{Does the ``binary-search reduction from detect to find'' (\S4 of the paper) preserve the empirical constant, or does it inflate the observed query count in a way that the big-$O$ statement hides?}\\
\emph{Basis:} our simulation tests only the \emph{detection} sub-problem (marked mass at the peak); the paper's headline number covers the \emph{find} problem, which pays a $\log$-search overhead.\\
\emph{Next steps:} chain a $\log N$-level binary search of the domain $Y$ using the walk as a subroutine; measure real total-query count for the find problem at $N\in\{8,10,12\}$ and see whether the $\log N$ multiplicative overhead is visible or absorbed.
\end{description}
